\section{The Conference in SQLite}
\label{sec:ch03-the-conference-in-sqlite}

Two packages, and the conference stops living in a static field.

\begin{minted}{text}
dotnet add package Microsoft.EntityFrameworkCore.Sqlite
dotnet add package HotChocolate.Data.EntityFramework
\end{minted}

\index{Entity Framework Core!SQLite provider}%
The first is the database provider and needs no introduction. The second is the
join between Hot Chocolate and EF Core, and it brings more with it than its own
name suggests: section~\ref{sec:ch03-asking-for-less} uses two types that
arrive through its dependencies rather than out of the package itself. The
project file gains both, and nothing else.

\begin{minted}[fontsize=\footnotesize, highlightlines={14,16}]{xml}
<Project Sdk="Microsoft.NET.Sdk.Web">

  <PropertyGroup>
    <TargetFramework>net10.0</TargetFramework>
    <Nullable>enable</Nullable>
    <ImplicitUsings>enable</ImplicitUsings>
    <GenerateDocumentationFile>true</GenerateDocumentationFile>
    <NoWarn>$(NoWarn);CS1591</NoWarn>
  </PropertyGroup>

  <ItemGroup>
    <PackageReference Include="HotChocolate.AspNetCore" Version="16.6.1" />
    <PackageReference Include="HotChocolate.AspNetCore.CommandLine" Version="16.6.1" />
    <PackageReference Include="HotChocolate.Data.EntityFramework" Version="16.6.1" />
    <PackageReference Include="HotChocolate.Types.Analyzers" Version="16.6.1" />
    <PackageReference Include="Microsoft.EntityFrameworkCore.Sqlite" Version="10.0.11" />
  </ItemGroup>

</Project>
\end{minted}

\index{DbContext@\code{DbContext}}%
\code{ConferenceContext.cs} is the whole of the mapping. Two tables, and the
seed data is the same \code{ConferenceData} class
chapter~\ref{ch:hot-chocolate-zero} printed, now read once at database creation
instead of on every request.

\begin{minted}{csharp}
using Microsoft.EntityFrameworkCore;

namespace Sessions;

/// <summary>
/// The conference program, in SQLite. Session carries SpeakerId as a
/// plain column and no navigation property, so EF Core never joins the
/// two tables on its own and every speaker lookup is a resolver's job.
/// </summary>
public sealed class ConferenceContext(
    DbContextOptions<ConferenceContext> options) : DbContext(options)
{
    public DbSet<Session> Sessions => Set<Session>();

    public DbSet<Speaker> Speakers => Set<Speaker>();

    protected override void OnModelCreating(ModelBuilder model)
    {
        // SQLite stores no date type of its own and refuses to sort a
        // DateTimeOffset. Keep the offset out of the column and put it
        // back on the way in; every session here is in conference time.
        model.Entity<Session>()
            .Property(s => s.StartsAt)
            .HasConversion(
                v => v.UtcDateTime,
                v => new DateTimeOffset(v, TimeSpan.Zero));

        model.Entity<Speaker>().HasData(ConferenceData.Speakers);
        model.Entity<Session>().HasData(ConferenceData.Sessions);
    }
}
\end{minted}

\index{SQLite!storage classes}%
\index{value converter}%
That converter is not decoration. SQLite has no date type of its own, and the
provider will not pretend otherwise: ask it to \code{ORDER BY} a
\code{DateTimeOffset} and it refuses to translate the query at all, naming the
type and telling you to convert it or to sort in memory instead. Since
\code{GetSessions} has ordered by \code{StartsAt} since
chapter~\ref{ch:hot-chocolate-zero}, that is the first thing to break. Storing
the instant and putting the offset back on the way out keeps the record, the
schema and every request already written exactly as they were. Sorting in
memory would have worked too, and would have meant reading every session on
every request forever, which is the opposite of where this chapter is going.

\index{Session@\code{Session}!has no navigation property}%
The other decision in that file is the one that is not there. \code{Session}
has a \code{SpeakerId} and no \code{Speaker} property, so EF Core sees an
\code{int} column and no relationship, and will never join the two tables on
its own initiative. A navigation property would have made this chapter's
problem disappear and taught you nothing, because
chapter~\ref{ch:router-n-plus-1} puts a network between these two tables and
there is no navigation property that spans it. The seam is easier to see if you
build it now, while both halves are still in the same process.

\index{statement log}%
\index{DbCommandInterceptor@\code{DbCommandInterceptor}}%
Next, the instrument. A database you cannot watch is a database you will guess
about, and the guesses in this particular area are famously wrong, so
\code{StatementLog.cs} prints every statement the service sends and numbers it.

\begin{minted}{csharp}
using System.Data.Common;
using Microsoft.EntityFrameworkCore.Diagnostics;

namespace Sessions;

/// <summary>
/// Prints every statement the service sends to SQLite, numbered, so a
/// request's cost can be read off the console. The counter is
/// process-wide and never reset: start the service, send one request,
/// and the last number is what that request cost.
/// </summary>
public sealed class StatementLog : DbCommandInterceptor
{
    private static int _count;

    public override InterceptionResult<DbDataReader> ReaderExecuting(
        DbCommand command,
        CommandEventData eventData,
        InterceptionResult<DbDataReader> result)
    {
        Print(command);
        return result;
    }

    public override ValueTask<InterceptionResult<DbDataReader>>
        ReaderExecutingAsync(
            DbCommand command,
            CommandEventData eventData,
            InterceptionResult<DbDataReader> result,
            CancellationToken cancellationToken = default)
    {
        Print(command);
        return ValueTask.FromResult(result);
    }

    private static void Print(DbCommand command)
    {
        var n = Interlocked.Increment(ref _count);
        Console.WriteLine($"-- statement {n}");
        Console.WriteLine(command.CommandText);
    }
}
\end{minted}

EF Core will log commands for you and I am not using that here, for a reason
worth stating once. Its own logging reports each command after it completes,
with the elapsed time attached, and this book prints no timings anywhere: a
duration is a fact about my laptop, and a statement count is a fact about the
program you typed. Counting on the way in also means the number is the same
whether the query returned in a millisecond or timed out.

\index{Program.cs@\code{Program.cs}!database registration}%
\code{Program.cs} carries the wiring. Everything
chapter~\ref{ch:hot-chocolate-zero} put there is untouched at the bottom of the
file; the new material is at the top.

\begin{minted}[highlightlines={3,4,8-24,29-31,35-41}]{csharp}
using ChilliCream.Nitro.App;
using HotChocolate.AspNetCore;
using Microsoft.EntityFrameworkCore;
using Sessions;

var builder = WebApplication.CreateBuilder(args);

// EF Core reports every command through the logging pipeline as well,
// at Information and with a duration attached. StatementLog below is
// this service's instrument, so silence the other one rather than have
// every statement printed twice.
builder.Logging.AddFilter(
    "Microsoft.EntityFrameworkCore.Database.Command",
    LogLevel.Warning);

builder.Services.AddDbContextFactory<ConferenceContext>(options =>
{
    options.UseSqlite("Data Source=conference.db");

    if (builder.Environment.IsDevelopment())
    {
        options.AddInterceptors(new StatementLog());
    }
});

builder
    .AddGraphQL()
    .AddSessionsTypes()
    .RegisterDbContextFactory<ConferenceContext>()
    .AddQueryContext()
    .AddPagingArguments();

var app = builder.Build();

using (var scope = app.Services.CreateScope())
{
    var factory = scope.ServiceProvider
        .GetRequiredService<IDbContextFactory<ConferenceContext>>();
    using var db = factory.CreateDbContext();
    db.Database.EnsureCreated();
}

app.MapGraphQL()
    .WithOptions(options =>
    {
        // Serve the Nitro build that shipped in the package. The
        // default, ServeMode.Latest, downloads the current one
        // from ChilliCream's CDN instead.
        options.Tool.ServeMode = ServeMode.Embedded;
    });

app.RunWithGraphQLCommands(args);
\end{minted}

\index{RegisterDbContextFactory@\code{RegisterDbContextFactory}}%
Two of those lines are worth stopping on. The \code{using} for the
\code{Sessions} namespace is needed because \code{Program.cs} is a top-level
program in no namespace, and this is the first file in the book that names a
type of yours rather than only calling generated methods. And
\code{AddDbContextFactory} does double duty: it teaches the container how to
build a context on demand, and it also registers the context itself as a
scoped service, which is what makes a \code{ConferenceContext} parameter on a
resolver resolvable at all.

\code{RegisterDbContextFactory} is the line I can tell you least about. It is
what the documentation pairs with a factory registration, and it is here
because I would ship it. Taking it out, though, leaves the exported schema,
every response and every statement count in this chapter unchanged, so I will
not tell you what it buys you.

\index{AddQueryContext@\code{AddQueryContext}}%
\index{AddPagingArguments@\code{AddPagingArguments}}%
\code{AddQueryContext} and \code{AddPagingArguments} belong to
section~\ref{sec:ch03-asking-for-less} and do nothing until something uses
them. They are here rather than there because
chapter~\ref{ch:hot-chocolate-zero} printed this file whole and
decision-by-decision reprints of a program get tiresome; this is the last time
\code{Program.cs} appears in the chapter.

\index{EnsureCreated@\code{EnsureCreated}}%
The block after \code{Build()} creates the database file and seeds it on first
run, and the call it makes is the blunt one. \code{EnsureCreated} builds the
schema if none is there and does nothing at all if a file already exists, so
it never migrates, and it will leave a stale database in place after you
change a record. Delete the file when that happens. Migrations are the
grown-up answer, and they generate files this book would then have to print,
for a two-table example that is rebuilt from nothing in a second.

\index{resolver!taking a DbContext}%
The two resolvers change shape. \code{Query.cs} first, taking the context as a
parameter and awaiting the result:

\begin{minted}{csharp}
using Microsoft.EntityFrameworkCore;

namespace Sessions;

[QueryType]
public static class Query
{
    /// <summary>
    /// Every session on the schedule, earliest first.
    /// </summary>
    public static async Task<List<Session>> GetSessionsAsync(
        ConferenceContext db,
        CancellationToken ct)
        => await db.Sessions.OrderBy(s => s.StartsAt).ToListAsync(ct);

    /// <summary>
    /// One session by its identifier, or null if no session has it.
    /// </summary>
    public static async Task<Session?> GetSessionByIdAsync(
        int id,
        ConferenceContext db,
        CancellationToken ct)
        => await db.Sessions.FirstOrDefaultAsync(s => s.Id == id, ct);
}
\end{minted}

\index{Parent@\code{[Parent]}!with a service beside it}%
Then \code{SessionType.cs}, the file this chapter is really about. The
translation from the in-memory version is mechanical: the linear scan over
\code{ConferenceData.Speakers} becomes a query, and the method gains the
context and a cancellation token.

\begin{minted}{csharp}
using Microsoft.EntityFrameworkCore;

namespace Sessions;

[ObjectType<Session>]
public static partial class SessionType
{
    /// <summary>
    /// The person giving this session.
    /// </summary>
    public static async Task<Speaker?> GetSpeakerAsync(
        [Parent] Session session,
        ConferenceContext db,
        CancellationToken ct)
        => await db.Speakers
            .FirstOrDefaultAsync(s => s.Id == session.SpeakerId, ct);
}
\end{minted}

Nothing about that is wrong. It is the obvious translation, it is what the
resolver did before, and it is what I would write first. Send
chapter~\ref{ch:hot-chocolate-zero}'s request and every response comes back
exactly as it did, which is what the move was supposed to achieve: the data
moved and the API did not.

\index{cost directive!attached to a resolver}%
Export the schema anyway. One line of it is different, and it is not a line you
wrote. Where chapter~\ref{ch:hot-chocolate-zero} exported \code{speaker:
Speaker}, the same field now reads:

\begin{graphqlsdl}
speaker: Speaker @cost(weight: "10")
\end{graphqlsdl}

Two things about the new method are enough to attach that weight, and either
does it alone: returning a \code{Task}, and taking a service parameter.
Compile chapter~\ref{ch:hot-chocolate-zero}'s method unchanged in this project
and the field exports bare again. Hot Chocolate is reading the signature and
concluding that this field might cost something, which is the same conclusion
this section has just reached by hand, and it is worth knowing before
chapter~\ref{ch:composition}, where a second tool reads these weights and does
not like what it finds.

Everything else about the service is what it was, and it is now more expensive
than you can see from outside it.
